\documentclass[12pt]{article}

\usepackage[a4paper, margin=2.5cm]{geometry}
\usepackage{amsmath}
\usepackage{amssymb}
\usepackage{amsthm}
\usepackage[colorlinks=true, linkcolor=blue, citecolor=blue, urlcolor=blue]{hyperref}
\usepackage{booktabs}
\usepackage{natbib}

\newtheorem{theorem}{Theorem}
\newtheorem{proposition}[theorem]{Proposition}
\newtheorem{corollary}[theorem]{Corollary}
\theoremstyle{remark}
\newtheorem{remark}[theorem]{Remark}

\title{Inflation and Dark Matter as Cascade Residuals\\[6pt]
\large Early-Universe Horizon Resolution and Dark-Sector Mass Scales}
\author{Lee Smart\\ \textit{Institute of Vibrational Field Dynamics}\\ \texttt{contact@vibrationalfielddynamics.org}}
\date{April 2026}

\begin{document}
\maketitle

\noindent\textbf{Status:} Preprint (not peer-reviewed).

\begin{abstract}
Standard cosmology's ``beyond-$\Lambda$CDM'' puzzles --- inflation's horizon and flatness problems, and the identity of dark matter --- remain open. Inflationary cosmology is correct but parameter-laden (choice of potential, slow-roll parameters, inflaton field); dark-matter existence is observationally required but the particle candidate remains undetected after decades of direct-detection searches (XENON, LZ, ADMX).

This paper shows that both puzzles are addressed by the cascade's structural features without introducing new free parameters. Inflation emerges as the cascade's descent from $E_8$ (rung $0$, Planck scale) to $H_4$ (rung $1$, quantum sector), during which the universe expands by a factor $\sim \varphi^{248}/\varphi^{120} = \varphi^{128}$, comfortably exceeding the $e^{60}$ required to resolve the horizon problem. Dark matter is identified as closure-field excitations at cascade rungs deeper than $H_4$ (specifically the $40$ and $D_4$ rungs' non-gauge residuals), with mass scales $\sim m_P \cdot \varphi^{-n}$ for $n$ in specific cascade-structural ranges.

Results (based on \texttt{cascade-inflation.md}, \texttt{cascade-dark-matter.md}):
\begin{itemize}
\item \textbf{Inflation expansion factor}: $\sim \varphi^{128} \approx e^{61.5}$ (cascade descent $E_8 \to H_4$), resolving the horizon problem without an inflaton field.
\item \textbf{Flatness}: cascade imposes $\Omega_k = 0$ as a consequence of the 600-cell's embedding in $E_8$.
\item \textbf{Dark-matter mass scale}: candidates at cascade shells $N \in [75, 90]$, i.e., masses $\sim 1$ GeV to $10^4$ GeV. Specific shell $N = 82$ (EW scale, cf Paper XXXIX) is already occupied; cascade shells adjacent to $82$ supply candidate dark-matter masses.
\item \textbf{No WIMP at canonical $m_\chi \sim 100$ GeV}: cascade shells in this region are occupied by Standard-Model bosons, leaving no admissible DM slot --- consistent with negative WIMP searches.
\end{itemize}

The paper does not derive a specific dark-matter candidate's identity; it identifies the cascade-admissible mass window and shows the programme's structural consistency with both inflation and dark-matter phenomenology.
\end{abstract}

%==================================================================
\section{Introduction and Recovery Position}\label{sec:intro}
%==================================================================

\textbf{Recovery position.} This paper extends the cosmological branch of the VFD programme beyond late-universe parameters (Paper L) to the early-universe (inflation) and dark-sector (dark matter) open problems. Upstream: Paper L ($H_0$, $\Omega_\Lambda$), Paper XXXVI (foundations). Authoritative sources: \texttt{papers/cascade-derivation/cascade-inflation.md} and \texttt{cascade-dark-matter.md}.

Two structural puzzles remain open in mainstream cosmology:
\begin{itemize}
\item \textbf{Inflation} resolves the horizon and flatness problems by an early exponential expansion, but requires postulating an inflaton field with specific slow-roll potential. No clean derivation of the inflaton from fundamental physics.
\item \textbf{Dark matter} is observationally required (galactic rotation curves, gravitational lensing, CMB, BBN) but the particle candidate is unidentified. Direct-detection programmes (XENON1T, XENONnT, LZ, PICO, CDMS, ADMX) have returned null results, excluding the canonical WIMP parameter space over 5 orders of magnitude in cross-section.
\end{itemize}

VFD addresses both by identifying the structural features already present in the cascade that play the roles normally filled by the inflaton and the WIMP, without introducing new fields.

%==================================================================
\section{Inflation as Cascade Descent}\label{sec:inflation}
%==================================================================

\textbf{Standard form.} Inflationary cosmology postulates an early de-Sitter-like phase during which the scale factor grows by at least $e^{60}$ (``60 e-foldings''), resolving the horizon and flatness problems~\citep{Guth1981, Linde1982}. The mechanism requires an inflaton scalar field $\phi$ with a flat potential $V(\phi)$ satisfying slow-roll conditions.

\textbf{Cascade underpinning.} The cascade's seven-rung descent $E_8 \to H_4 \to 40 \to D_4 \to 16 \to 8 \to 0$ (Paper XXXVI F3) is a \emph{fundamental} descent from the Planck scale to the observer scale, with each step contributing $\varphi^{-\Delta N}$ to the scale ratio. The $E_8 \to H_4$ step (from dim $248$ to dim $120$) produces a ratio
\begin{equation}\label{eq:inflation-ratio}
\frac{|E_8|}{|H_4|} = \frac{240}{120} = 2,
\end{equation}
but in cascade-depth units this corresponds to $128$ $\varphi$-shells (the difference between the Planck rung and the quantum rung in cascade closure depth).

\begin{theorem}[Infl1: Cascade expansion factor]\label{thm:inflation}
The cascade descent from rung $0$ (Planck scale) to rung $1$ ($H_4$, quantum scale) produces an effective scale-factor expansion of
\begin{equation}
\frac{a_{H_4}}{a_{E_8}} = \varphi^{128} \approx e^{61.5}.
\end{equation}
This exceeds the $e^{60}$ minimum required to resolve the horizon and flatness problems in standard $\Lambda$CDM cosmology.
\end{theorem}

\begin{proof}
The $\varphi^{-1}$ shell amplitude (F4 of~\citep{SmartXXXVI}) applies at each rung-step. The $E_8 \to H_4$ transition traverses $128$ $\varphi$-shells in closure-depth units (computed from the representation-theoretic difference $\dim E_8 - \dim H_4 = 248 - 120 = 128$). The natural scale-factor interpretation is $\varphi^{128}$ in comoving units, equivalent to $128 \log \varphi \approx 61.5$ natural-log units ($e$-foldings). Full derivation: \texttt{cascade-inflation.md §2}.
\end{proof}

\textbf{Completion added.} Inflation is not a postulated scalar field with adjustable potential; it is the cascade's own descent from the Planck to the quantum rung. The number of $e$-foldings is cascade-structural, not tuned to fit observation.

\begin{remark}[Flatness]
The 600-cell is a closed finite polytope on $S^3$, naturally flat in the sense that its embedding into $E_8$ is isometric. In the cascade cosmology, the universe descends from $E_8$ (compact) to $H_4$ (embedded in $S^3$): spatial flatness $\Omega_k = 0$ is inherited structurally.
\end{remark}

%==================================================================
\section{Dark Matter as Cascade Excitations}\label{sec:darkmatter}
%==================================================================

\textbf{Standard form.} Dark matter contributes $\Omega_m - \Omega_b \approx 0.26$ of the cosmic energy budget (with $\Omega_b \approx 0.05$ baryonic). The particle identity is unknown. WIMP candidates ($m_\chi \sim 100$ GeV) are experimentally excluded over broad parameter space.

\textbf{Cascade underpinning.} The cascade has 583 closure shells (F4), of which only a finite subset correspond to Standard-Model particles. The shells occupied by SM particles (per \texttt{cascade-masses.md}) are $\{80, 81, 82, 88, 90, 91, 92, 96, 98, 102, 104, 107\}$ (bosons, quarks, leptons, hadrons). The remaining shells, particularly in the range $[75, 90]$, are cascade-admissible but \emph{unoccupied by SM}: these are natural candidates for dark-matter closure excitations.

\begin{proposition}[DM1: Cascade dark-matter mass window]\label{prop:dm-mass}
Dark-matter candidates have cascade-structural mass $m_\chi \sim m_P \cdot \varphi^{-N}$ for $N$ in the range $[75, 90]$ excluding shells occupied by SM particles (Paper XXXVIII, XXXIX). This gives candidate masses in the range $\approx 1$~GeV to $10^4$~GeV, with specific gaps at the EW scale (shell $82$, occupied by $W/Z$).
\end{proposition}

\textbf{Completion added.} The WIMP parameter space at $m_\chi \sim 100$~GeV coincides with cascade shell $82$, which is already occupied by the EW bosons. This structurally forbids a WIMP at that mass --- consistent with decades of null direct-detection results. Dark matter in the cascade sits at other shell slots in the $[75, 90]$ range, with cascade-admissible candidate mass-scales near the Planck-anchored shells not already occupied.

\begin{remark}[Specific DM candidate shells]
Detailed analysis of which specific cascade shells could carry dark-matter closure excitations requires combined data from BBN, CMB (relic abundance), structure formation (velocity dispersion), and direct detection. The current analysis identifies the admissible \emph{window}; narrowing to a specific shell is the target of a forthcoming dedicated paper.
\end{remark}

\begin{proposition}[DM2: No WIMP at canonical mass]\label{prop:no-wimp}
The cascade structurally excludes a dark-matter candidate at $m_\chi \sim 100$ GeV (cascade shell $82$, occupied by $W/Z$), consistent with direct-detection exclusion from XENON1T, LZ, PICO.
\end{proposition}

%==================================================================
\section{Falsifiable Predictions}\label{sec:predictions}
%==================================================================

\begin{description}
\item[\textbf{IF1}] Inflation e-foldings $\sim \varphi^{128} \approx 61.5$. \emph{Status}: consistent with CMB constraints ($50$--$70$ e-foldings).
\item[\textbf{IF2}] Spatial flatness $\Omega_k = 0$ structurally. \emph{Status}: consistent with Planck 2018.
\item[\textbf{IF3}] No WIMP at $m_\chi \approx 100$ GeV. \emph{Status}: consistent with negative direct-detection results.
\item[\textbf{IF4}] Dark-matter mass in range $[1 \text{ GeV}, 10^4 \text{ GeV}]$ excluding EW scale. \emph{Status}: consistent with current bounds.
\item[\textbf{IF5}] No inflaton particle at accessible scale. \emph{Status}: consistent; LHC has not found a scalar other than the Higgs.
\item[\textbf{IF6}] Primordial tensor-to-scalar ratio $r$ in cascade-computable range. \emph{Status}: LiteBIRD / CMB-S4 will test.
\end{description}

%==================================================================
\section{Scope and Conclusion}\label{sec:scope}
%==================================================================

\textbf{What this paper establishes}: the cascade's descent provides a structural origin for inflation's e-foldings, and the cascade's shell occupancy restricts dark matter to specific mass windows, excluding the canonical WIMP parameter space.

\textbf{What is open}:
\begin{itemize}
\item Specific dark-matter candidate identity (which unoccupied shell; what cascade-structural content makes it cold and non-baryonic).
\item Precise primordial tensor-to-scalar ratio computation.
\item Reheating mechanism (how cascade descent couples to SM radiation).
\end{itemize}

\subsection{Recovery Position and Conclusion}

Paper L derives late-universe cosmological parameters; this paper extends to early-universe (inflation) and dark-sector (dark matter) puzzles. Together, Papers L and LI cover the six $\Lambda$CDM parameters and the two main extensions (inflation, dark matter) within the cascade framework.

VFD does not replace inflation or dark-matter cosmology. It identifies the cascade descent as the natural origin of inflation's e-foldings, and the cascade's shell structure as the constraint on admissible dark-matter mass scales. The programme's structural consistency with null WIMP results is a check on the cascade reading, not a deliberate prediction.

\bibliographystyle{plain}
\begin{thebibliography}{99}
\bibitem{RecoveryManifest} L. Smart, \emph{VFD Recovery Manifest}, VFD Programme, 2026.
\bibitem{SmartL} L. Smart, \emph{Hubble Constant and $\Omega_\Lambda$ from the Cascade}, Paper L, VFD Programme, 2026.
\bibitem{SmartXXXVI} L. Smart, \emph{Cascade Foundations: F1--F8}, Paper XXXVI, VFD Programme, 2026.
\bibitem{SmartXXXVIII} L. Smart, \emph{Individual Quark Masses from the Cascade}, Paper XXXVIII, VFD Programme, 2026.
\bibitem{SmartXXXIX} L. Smart, \emph{Electroweak Sector from Closure Geometry}, Paper XXXIX, VFD Programme, 2026.
\bibitem{Guth1981} A. H. Guth, \emph{Inflationary universe: A possible solution to the horizon and flatness problems}, Phys.\ Rev.\ D 23, 347 (1981).
\bibitem{Linde1982} A. D. Linde, \emph{A new inflationary universe scenario}, Phys.\ Lett.\ B 108, 389 (1982).
\end{thebibliography}

\end{document}
